\chapter{Risk Questionnaire Instrument}
\label{appendix:risk-questionnaire}

This appendix documents the production questionnaire used by the platform to estimate user risk tolerance and map users to recommendation profiles.

\textbf{Questionnaire Overview}

\begin{itemize}[leftmargin=*]
	\item Title: ETF Goal \& Risk Quick Check
	\item Total questions: 8 (multiple-choice)
	\item Options per question: 4 (A--D)
	\item Scoring method: sum of option points across answered questions
	\item Output profiles: Conservative, Balanced, Aggressive
\end{itemize}

\textbf{Scoring Notes}

\begin{itemize}[leftmargin=*]
	\item Backend score logic sums points from selected options. If an answer is missing, that question is skipped by API logic.
	\item In normal frontend flow, users must answer all 8 questions before submission.
	\item Backend percentage uses configured \texttt{max\_score = 720} from quiz configuration. Based on the current options, the effective highest attainable raw total is 715.
	\item Profile thresholds are configuration-driven in backend quiz data and can be adjusted without changing API contracts.
	\item The questionnaire supports recommendation personalization and does not constitute regulated investment advice.
\end{itemize}

\section{Question Set (8 Questions)}

\begin{longtable}{|>{\raggedright\arraybackslash\bfseries}p{1.1cm}|p{9.1cm}|p{3.2cm}|}
\hline
\rowcolor{gray!15} QID & Question Text & Answer Options (Points) \\ \hline
Q1 & When do you expect to need this money? & A: Within 3 months (5), B: 3--12 months (20), C: 1--3 years (55), D: More than 3 years (90) \\ \hline
Q2 & What is your primary purpose for this investment? & A: Protect my principal with minimal fluctuation (10), B: Beat bank savings rate with low risk (25), C: Balanced growth over time (60), D: Maximize long-term growth (90) \\ \hline
Q3 & If your ETF position drops 8\% in one month, what would you likely do? & A: Sell immediately to avoid further loss (10), B: Reduce part of my position (35), C: Hold and wait (65), D: Add more if fundamentals still look good (90) \\ \hline
Q4 & How much downside can you accept before this investment causes real stress? & A: About 2\% or less (10), B: Around 5\% (35), C: Around 10\% (65), D: 15\% or more (90) \\ \hline
Q5 & Do you already have an emergency fund (at least 3--6 months of expenses)? & A: No (10), B: Partially (1--3 months) (35), C: Yes (3--6 months) (65), D: Yes (more than 6 months) (85) \\ \hline
Q6 & Which best describes this money amount? & A: Critical cash I may need soon (10), B: Important, but not urgent (35), C: Dedicated investment capital (65), D: High-risk growth capital I can leave untouched (90) \\ \hline
Q7 & How familiar are you with ETFs and market behavior? & A: Beginner (10), B: Basic understanding (35), C: Comfortable selecting broad ETFs (65), D: Advanced and comfortable with volatility (90) \\ \hline
Q8 & What return-vs-stability tradeoff fits you best? & A: I prefer stable value even if returns are lower (10), B: I want modest outperformance vs savings with low swings (30), C: I accept moderate fluctuations for better growth (65), D: I prioritize growth and can tolerate larger ups and downs (90) \\ \hline
\end{longtable}

\section{Profile Thresholds}

\begin{table}[H]
\centering
\begin{tabular}{lccp{7.2cm}}
\toprule
\textbf{Profile} & \textbf{Min Score} & \textbf{Max Score} & \textbf{Description} \\
\midrule
Conservative & 0 & 260 & Low-risk investor prioritizing capital preservation and stable income \\
Balanced & 261 & 500 & Moderate-risk investor balancing growth with stability \\
Aggressive & 501 & 720 & High-risk investor prioritizing capital growth over stability \\
\bottomrule
\end{tabular}
\caption{Risk Profile Mapping for the 8-Question Questionnaire}
\label{tab:appendix-risk-thresholds}
\end{table}
